%============================================================
%  ch13_lindblad.tex
%  Chapter 13: The Dirac--Frenkel Principle for Open Quantum
%  Systems
%
%  Tensor Formats in Banach Spaces
%  Antonio Falcó, CEU Universities
%
%  Based on: Acevedo--Falcó (2026), FHN2019
%  Estimated length: ~30 pp.
%============================================================

\chapter{The Dirac--Frenkel Principle for Open Quantum Systems}
\label{chap:lindblad}

Chapter~\ref{chap:DF} established the Dirac--Frenkel variational
principle for abstract Cauchy problems on tensor Banach manifolds.
Chapter~\ref{chap:quantum_states} identified the fixed-rank stratum
$\Drank{r}{H}$ as a smooth Banach manifold within the trace-class
Banach space $\cTsa{1}(H)$, inheriting all the geometric properties
developed in Chapters~\ref{chap:tucker_manifold}--\ref{chap:TB_manifold}.
This chapter combines these two lines to obtain the Dirac--Frenkel
principle for quantum master equations on $\cTsa{1}(H)$: the Lindblad
(GKSL) equation for open quantum systems.

The Lindblad equation describes the evolution of the state $\rho(t)
\in \Dstate{H}$ of a quantum system coupled to an environment. Its
generator $\Lind$ maps $\cTsa{1}(H)$ to itself and preserves
$\Dstate{H}$. Since exact simulation of $\rho(t)$ in
infinite-dimensional systems is generally intractable, one seeks a
reduced dynamics on the finite-rank manifold $\Drank{r}{H}$,
obtained by projecting the Lindblad vector field onto the tangent space
of $\Drank{r}{H}$.

Section~\ref{sec:lind:equation} formulates the Lindblad equation as an
ODE on $\cTsa{1}(H)$ and identifies its vector field.
Section~\ref{sec:lind:tangent} develops the tangent space of $\Drank{r}{H}$
via minimal subspaces and establishes the bounded projector.
Section~\ref{sec:lind:df} states and proves the main well-posedness theorem.
Section~\ref{sec:lind:preservation} analyses which structural properties
(trace, rank, positivity) are preserved by the projected dynamics.
Section~\ref{sec:lind:kinematic} derives the kinematic equations in
factorised coordinates, recovering the equations of Chapter~\ref{chap:DF}
as a special case.
Section~\ref{sec:lind:history} contains historical notes.

%------------------------------------------------------------
\section{The Lindblad equation on $\cTsa{1}(H)$}
\label{sec:lind:equation}
%------------------------------------------------------------

\subsection*{The GKSL generator}

The most general Markovian master equation for the state $\rho(t)
\in \Dstate{H}$ of an open quantum system is the
\emph{Gorini--Kossakowski--Sudarshan--Lindblad} (GKSL) equation:
\begin{equation}
\label{eq:lind:GKSL}
    \dot\rho(t) = \Lind(\rho(t))
    \coloneqq -i[H, \rho(t)]
    + \sum_{j \in J} \Bigl(
    L_j \rho(t) L_j^*
    - \tfrac{1}{2}\{L_j^* L_j, \rho(t)\}
    \Bigr),
\end{equation}
where:
\begin{itemize}
\item $H = H^* \in \cL(H)$ is the \emph{Hamiltonian} (self-adjoint,
    bounded);
\item $\{L_j\}_{j \in J} \subset \cL(H)$ are the \emph{Lindblad
    operators} (noise operators), with $J$ a finite or countable index set;
\item $[H, \rho] = H\rho - \rho H$ is the commutator and
    $\{A, B\} = AB + BA$ is the anticommutator.
\end{itemize}
When $J$ is finite and all $L_j \in \cL(H)$ are bounded, $\Lind$ is a
bounded operator on $\cT{1}(H)$. The solution to~\eqref{eq:lind:GKSL}
is a strongly continuous one-parameter semigroup $(T_t)_{t \geq 0}$ on
$\cT{1}(H)$ that is \emph{completely positive and trace preserving}
(CPTP): $T_t(\Dstate{H}) \subset \Dstate{H}$ for all $t \geq 0$.

\begin{remark}[The Lindblad equation as a tensor Banach ODE]
\label{rem:lind:tensor_ODE}
The equation~\eqref{eq:lind:GKSL} is a special case of the abstract
Cauchy problem~\eqref{eq:DF:ode} of Chapter~\ref{chap:DF}, with ambient
Banach space $\VDnorm = \cTsa{1}(H)$ and vector field
$\bF(t, \rho) = \Lind(\rho)$ (time-independent in the homogeneous case).
The trace norm $\|\cdot\|_1$ plays the role of $\|\cdot\|_D$, and the
identification $\cTsa{1}(H) \cong H \otimes_\pi H^*$ (Remark~\ref{rem:app:trace_projective})
makes $(\cTsa{1}(H), \|\cdot\|_1)$ a Banach tensor space with $d=2$,
$V_1 = H$, $V_2 = H^*$, satisfying condition~\condINJ\ with the
projective norm.
\end{remark}

\begin{assumption}[Regularity of $\Lind$]
\label{ass:lind:regular}
Throughout this chapter we assume that:
\begin{enumerate}
\item[\textup{(R1)}] $H \in \cL(H)$ is bounded and $J$ is finite,
    so $\Lind: \cTsa{1}(H) \to \cTsa{1}(H)$ is a bounded linear operator.
\item[\textup{(R2)}] The CPTP semigroup $(T_t)_{t \geq 0}$ generated
    by $\Lind$ satisfies $T_t(\Dstate{H}) \subset \Dstate{H}$.
\item[\textup{(R3)}] For each $\rho_0 \in \Drank{r}{H}$, the restriction
    $\Lind|_{\Drank{r}{H}}$ is locally Lipschitz in $\|\cdot\|_1$.
\end{enumerate}
\end{assumption}

Under Assumption~\ref{ass:lind:regular}, the full Lindblad equation has
a unique global solution $\rho: [0,\infty) \to \Dstate{H}$.

\begin{remark}[Unbounded generators]
\label{rem:lind:unbounded}
Assumption~\ref{ass:lind:regular}(R1) excludes unbounded Lindblad
operators (e.g., position or momentum operators in infinite dimensions).
The extension to unbounded generators requires domain-adapted fixed-rank
strata and is the subject of Open Problem~\ref{prob:lindblad_unbounded}.
\end{remark}

%------------------------------------------------------------
\section{Tangent space of $\Drank{r}{H}$ via minimal subspaces}
\label{sec:lind:tangent}
%------------------------------------------------------------

The tangent space characterisation of Theorem~\ref{thm:qs:manifold}(iii)
takes the following explicit form, connecting with the Tucker tangent
space of Chapter~\ref{chap:tucker_manifold}.

\begin{proposition}[Tangent space decomposition]
\label{prop:lind:tangent}
Let $\rho \in \Drank{r}{H}$, $S = \supp(\rho)$, $P = \PS{S}$.
With respect to the orthogonal decomposition $H = S \oplus S^\perp$,
every $\dot\rho \in T_\rho \Drank{r}{H}$ has a unique block form
\begin{equation}
\label{eq:lind:tangent_block}
    \dot\rho =
    \begin{pmatrix} C & B^* \\ B & 0 \end{pmatrix}
    \in \cTsa{1}(H),
    \quad C \in \cTsa{1}(S),\; \Tr(C) = 0,\;
    B \in \cT{1}(S, S^\perp),
\end{equation}
where $\cT{1}(S, S^\perp)$ denotes the trace-class operators from
$S$ to $S^\perp$ (automatically finite-rank since $\dim S = r$).
The decomposition $T_\rho \Drank{r}{H} \cong \{C \in \cTsa{1}(S) :
\Tr(C) = 0\} \oplus \cT{1}(S, S^\perp)$ identifies:
\begin{itemize}
\item The \emph{core direction} $C$: variation of $\rho$ within fixed
    support $S$ — the quantum analogue of $\dot{C}^{(D)}$ in~\eqref{eq:DF:tangent_decomp}.
\item The \emph{support direction} $B$: variation of the support
    subspace $S$ — the quantum analogue of $\dot{u}^{(\alpha)}_{i_\alpha}
    \in \Wmin{\alpha}{\bv_r}$ in~\eqref{eq:DF:tangent_decomp}.
\end{itemize}
\end{proposition}

\begin{proof}
This is Lemma~3.1 and Corollary~3.2 of~\cite{AcevedoFalco2026},
proved by differentiating the chart curve $(A(t), \sigma(t))
\mapsto V_{A(t)}\sigma(t)V_{A(t)}^*$ at $t=0$ in local chart
coordinates~\eqref{eq:qs:chart}. The block structure follows from
the orthogonality of $H = S \oplus S^\perp$ and the chain rule,
exactly as in the proof of Proposition~\ref{prop:tangent_space}.
\end{proof}

\begin{remark}[The trace constraint]
\label{rem:lind:trace_constraint}
The condition $\Tr(\dot\rho) = 0$ in~\eqref{eq:lind:tangent_block}
is the infinitesimal form of $\Tr(\rho(t)) = 1$. It corresponds to
the fact that $\Drank{r}{H}$ lies in the affine hyperplane
$\{\Tr = 1\} \subset \cTsa{1}(H)$, which is an additional constraint
not present in the Tucker manifold $\Tucker{\fr}{\VD} \subset \VDnorm$.
In the kinematic equations of Section~\ref{sec:lind:kinematic},
this constraint is enforced automatically by the trace-removal term
in the projector $\hat{P}_\rho$ (Proposition~\ref{prop:qs:projector}).
\end{remark}

%------------------------------------------------------------
\section{The Dirac--Frenkel principle for Lindblad dynamics}
\label{sec:lind:df}
%------------------------------------------------------------

\begin{definition}[Dirac--Frenkel projected Lindblad dynamics]
\label{def:lind:df}
The \emph{Dirac--Frenkel projected Lindblad dynamics} on $\Drank{r}{H}$
is the ODE
\begin{equation}
\label{eq:lind:df}
    \dot\rho(t) = \hat{P}_{\rho(t)}\bigl(\Lind(\rho(t))\bigr),
    \qquad \rho(0) = \rho_0 \in \Drank{r}{H},
\end{equation}
where $\hat{P}_\rho: \cTsa{1}(H) \to T_\rho \Drank{r}{H}$ is the
bounded projector of Proposition~\ref{prop:qs:projector}.
\end{definition}

The equation~\eqref{eq:lind:df} is the quantum specialisation of the
Dirac--Frenkel equation~\eqref{eq:DF_tucker} of Chapter~\ref{chap:DF},
with $\Tucker{\fr}{\VD}$ replaced by $\Drank{r}{H}$, $\VDnorm$ by
$\cTsa{1}(H)$, $\bF$ by $\Lind$, and $\Pmetric{\bv_r}$ by $\hat{P}_\rho$.

\begin{theorem}[Local well-posedness,
{\cite[Theorem~4.6]{AcevedoFalco2026}}]
\label{thm:lind:wellposed}
Under Assumption~\ref{ass:lind:regular}, for every
$\rho_0 \in \Drank{r}{H}$ there exist $T > 0$ and a unique
$\mathcal{C}^1$-curve
\begin{equation*}
    \rho: [0, T) \longrightarrow \Drank{r}{H}
\end{equation*}
satisfying~\eqref{eq:lind:df} with $\rho(0) = \rho_0$.
The solution depends continuously on $\rho_0$ in the trace norm.
\end{theorem}

\begin{proof}
We verify the three conditions of Theorem~\ref{thm:app:ode} (Picard--Lindelöf
on Banach manifolds, Appendix~\ref{app:manifolds}).

\medskip
\noindent\textbf{Step~1: $\Drank{r}{H}$ is a $\mathcal{C}^\infty$ Banach manifold.}

This is Theorem~\ref{thm:qs:manifold}. The model space for a chart at
$\rho_0$ is $\cL(\CC^r, S_0^\perp) \times \cTsa{1}(\CC^r)$, a Banach
space (infinite-dimensional in the first factor, finite-dimensional in
the second).

\medskip
\noindent\textbf{Step~2: The reduced vector field is locally Lipschitz.}

In the chart $\Phi_{S_0}$ at $\rho_0$, the reduced vector field
in coordinates is
\begin{equation*}
    F(A, \sigma)
    \coloneqq d\Phi_{S_0}^{-1}\bigl(\Phi_{S_0}(A,\sigma)\bigr)
    \bigl[\hat P_{\rho(A,\sigma)}\bigl(\Lind(\rho(A,\sigma))\bigr)\bigr],
\end{equation*}
where $\rho(A,\sigma) = V_A \sigma V_A^*$.

We verify local Lipschitz continuity in three sub-steps:

\emph{(a) $\Lind(\rho)$ is Lipschitz in $\|\cdot\|_1$.}
By Assumption~\ref{ass:lind:regular}(R1), $\Lind$ is bounded on
$\cT{1}(H)$: $\|\Lind(\rho) - \Lind(\rho')\|_1 \leq \|\Lind\|_{\cL(\cT{1})}
\|\rho - \rho'\|_1$.

\emph{(b) $\rho \mapsto \hat P_\rho$ is locally Lipschitz.}
Writing $\hat P_\rho(X) = P_{\supp(\rho)} X P_{\supp(\rho)} + \ldots$,
the dependence on $\rho$ enters only through the orthogonal projection
$P = P_{\supp(\rho)}$. In the chart, $\supp(\rho(A,\sigma)) =
\mathrm{Im}(V_A)$ and $P = V_A V_A^*$. The map $A \mapsto V_A V_A^*$
is $\mathcal{C}^\infty$ by Step~1 (Theorem~\ref{thm:qs:manifold}(i)),
hence Lipschitz on bounded sets.

\emph{(c) The chart map is Lipschitz.}
$d\Phi_{S_0}^{-1}$ is bounded on bounded sets since $\Phi_{S_0}$ is a
$\mathcal{C}^\infty$ diffeomorphism with bounded differential.

Combining (a)--(c) by the chain rule: $F$ is locally Lipschitz on
a neighbourhood of $(0, V_0^*\rho_0 V_0)$ in $\cL(\CC^r, S_0^\perp)
\times \cTsa{1}(\CC^r)$.

\medskip
\noindent\textbf{Step~3: Local existence and uniqueness by Picard--Lindelöf.}

By Theorem~\ref{thm:app:picard} (Picard--Lindelöf in Banach spaces),
applied to the coordinate ODE $\dot z(t) = F(z(t))$ in the Banach space
$\cL(\CC^r, S_0^\perp) \times \cTsa{1}(\CC^r)$, there exist $T > 0$ and
a unique $\mathcal{C}^1$ solution $z: [0,T) \to \cL(\CC^r, S_0^\perp)
\times \cTsa{1}(\CC^r)$ with $z(0) = (0, V_0^*\rho_0 V_0)$.
Pushing forward by $\Phi_{S_0}$, we obtain $\rho(t) = \Phi_{S_0}(z(t))$,
a $\mathcal{C}^1$ curve in $\Drank{r}{H}$ satisfying~\eqref{eq:lind:df}.

\medskip
\noindent\textbf{Step~4: Continuous dependence on $\rho_0$.}

Continuous dependence of $z(t)$ on $z(0)$ follows from Theorem~\ref{thm:app:picard}.
Transported through the $\mathcal{C}^\infty$ chart map $\Phi_{S_0}$,
this gives continuous dependence of $\rho(t)$ on $\rho_0$ in the
trace norm.
\end{proof}

\begin{remark}[Choice of projector]
\label{rem:lind:projector_choice}
The projector $\hat{P}_\rho$ in Definition~\ref{def:lind:df} is the
explicit block projector of Proposition~\ref{prop:qs:projector}. Other
choices of bounded projector $\pi_\rho: \cTsa{1}(H) \to T_\rho
\Drank{r}{H}$ lead to equivalent well-posedness results; the choice
affects the trajectory $\rho(t)$ but not the structural preservation
properties. In particular, in the Hilbert--Schmidt setting one may use
the orthogonal projection in $\cT{2}(H)$, which gives the TDVP
(time-dependent variational principle) projection of~\cite{LeBrisRouchon2013}.
\end{remark}

%------------------------------------------------------------
\section{Structure preservation}
\label{sec:lind:preservation}
%------------------------------------------------------------

\begin{proposition}[Trace, rank, and positivity]
\label{prop:lind:preservation}
Let $\rho: [0,T) \to \Drank{r}{H}$ be the solution of~\eqref{eq:lind:df}.
Then for all $t \in [0, T)$:
\begin{enumerate}
\item[\textup{(i)}] \emph{Trace preservation:}
    $\Tr(\rho(t)) = 1$.
\item[\textup{(ii)}] \emph{Rank preservation:}
    $\mathrm{rank}(\rho(t)) = r$.
\item[\textup{(iii)}] \emph{Positivity:}
    $\rho(t) \geq 0$.
\end{enumerate}
\end{proposition}

\begin{proof}
\textup{(i)}: Differentiating, $\frac{d}{dt}\Tr(\rho(t)) =
\Tr(\dot\rho(t)) = \Tr(\hat{P}_{\rho(t)}(\Lind(\rho(t)))) = 0$,
since $\hat{P}_\rho(X) \in T_\rho \Drank{r}{H}$ satisfies
$\Tr(\hat{P}_\rho(X)) = 0$ by~\eqref{eq:lind:tangent_block}.

\textup{(ii)} and \textup{(iii)}: The solution $\rho(t)$ is a curve
in $\Drank{r}{H}$ by Theorem~\ref{thm:lind:wellposed}. Since
$\Drank{r}{H} \subset \Dstate{H}$, every point satisfies
$\rho(t) \geq 0$ and $\mathrm{rank}(\rho(t)) = r$.
\end{proof}

\begin{remark}[Complete positivity is not guaranteed]
\label{rem:lind:cptp}
Proposition~\ref{prop:lind:preservation} guarantees that $\rho(t)
\in \Drank{r}{H} \subset \Dstate{H}$ (positive, trace one, rank $r$)
along the trajectory. It does \emph{not} guarantee that the reduced
flow $\rho_0 \mapsto \rho(t)$ is a CPTP map (quantum channel).
The projected dynamics is a nonlinear ODE on $\Drank{r}{H}$; the
linearity and complete positivity of the full Lindblad semigroup
$(T_t)_{t \geq 0}$ are not inherited by the projected flow.
Constructing a CPTP-preserving rank-$r$ approximation requires
additional structure, such as Kraus-form parametrisations;
see Open Problem~\ref{prob:lindblad_cptp}.
\end{remark}

%------------------------------------------------------------
\section{Kinematic equations in factorised coordinates}
\label{sec:lind:kinematic}
%------------------------------------------------------------

\subsection*{Stiefel-core parametrisation}

Every $\rho \in \Drank{r}{H}$ admits a factorisation
\begin{equation}
\label{eq:lind:stiefel_core}
    \rho = V \sigma V^*,
    \quad V \in \Stiefel{r}{H},\;
    \sigma \in \Dpos{\CC^r},
\end{equation}
where $\Stiefel{r}{H} = \{V \in \cL(\CC^r, H) : V^*V = I_{\CC^r}\}$
is the Stiefel manifold of isometries. The parametrisation is unique up
to the unitary gauge action $(V, \sigma) \sim (VU, U^*\sigma U)$ for
$U \in U(r)$.

Differentiating~\eqref{eq:lind:stiefel_core}:
\begin{equation}
\label{eq:lind:diff_stiefel}
    \dot\rho = \dot{V} \sigma V^* + V \dot\sigma V^* + V \sigma \dot{V}^*.
\end{equation}
The Stiefel constraint $V^*V = I$ differentiates to give $V^*\dot{V} +
\dot{V}^*V = 0$, so $V^*\dot{V}$ is skew-Hermitian. Imposing the
\emph{horizontal (parallel transport) gauge}
\begin{equation}
\label{eq:lind:gauge}
    V^* \dot{V} = 0
    \quad\Longleftrightarrow\quad
    \dot{V} = (I - VV^*) \dot{V},
\end{equation}
the variation $\dot{V}$ lies in the complement of $\mathrm{Im}(V) = S$,
i.e., $\dot{V}(\CC^r) \subset S^\perp$.

\subsection*{Derivation of the kinematic system}

Under the gauge condition~\eqref{eq:lind:gauge}, projecting the
Lindblad vector field $\Lind(\rho)$ onto $T_\rho \Drank{r}{H}$ via
$\hat{P}_\rho$ gives, using the block formula~\eqref{eq:qs:projector}:

\medskip
\noindent\textbf{Equation for the core $\sigma(t)$.}
Multiplying $\dot\rho = \hat{P}_\rho(\Lind(\rho))$ on the left by $V^*$
and on the right by $V$, and using $V^*\dot{V} = 0$:
\begin{equation}
\label{eq:lind:core_eq}
    \dot\sigma(t) = V(t)^*\, \hat{P}_{\rho(t)}\bigl(\Lind(\rho(t))\bigr)\, V(t).
\end{equation}
Since $\hat{P}_\rho(X) = PXP + \text{off-diagonal} - \frac{\Tr(PXP)}{r}P$
and $P = VV^*$, this simplifies to
\begin{equation}
\label{eq:lind:core_explicit}
    \dot\sigma(t)
    = V^* \Lind(\rho) V
    - \frac{\Tr(V^* \Lind(\rho) V)}{r}\, I_{\CC^r}
    = V^* \Lind(\rho) V,
\end{equation}
where the last equality uses $\Tr(\Lind(\rho)) = 0$ (trace preservation
of the full Lindblad semigroup), which gives $\Tr(V^*\Lind(\rho)V) =
\Tr(VV^*\Lind(\rho)) = \Tr(P\Lind(\rho)) = 0$ when the $(S^\perp,S^\perp)$
block of $\Lind(\rho)$ vanishes on the tangent space.

\medskip
\noindent\textbf{Equation for the frame $V(t)$.}
Applying $(I - VV^*)$ on the left and $V\sigma^{-1}$ on the right to
$\dot\rho = \hat{P}_\rho(\Lind(\rho))$:
\begin{equation}
\label{eq:lind:frame_eq}
    \dot{V}(t)
    = (I - V(t)V(t)^*)\,
    \hat{P}_{\rho(t)}\bigl(\Lind(\rho(t))\bigr)\,
    V(t)\, \sigma(t)^{-1}.
\end{equation}
This is well-defined since $\sigma(t) \in \Dpos{\CC^r}$ is invertible
(positive definite) along the trajectory.

\begin{remark}[Equations~\eqref{eq:lind:core_explicit}--\eqref{eq:lind:frame_eq}
as a special case of the MCTDH system]
\label{rem:lind:mctdh_special}
Equations~\eqref{eq:lind:core_explicit}--\eqref{eq:lind:frame_eq} are
the quantum analogue of the MCTDH kinematic
system~\eqref{eq:DF:MCTDH_core}--\eqref{eq:DF:MCTDH_basis} of
Chapter~\ref{chap:DF}. Precisely:
\begin{itemize}
\item $\sigma(t) \in \Dpos{\CC^r}$ plays the role of the core tensor
    $C^{(D)}(t)$: the finite-dimensional degrees of freedom within the
    fixed support.
\item $V(t) \in \Stiefel{r}{H}$ plays the role of the basis
    $\{u^{(\alpha)}_{i_\alpha}(t)\}$ of $\Umin{\alpha}{\bv_r(t)}$:
    the isometric frame of the support subspace.
\item The gauge condition $V^*\dot{V} = 0$ is the quantum analogue of
    the tangential constraint~\eqref{eq:DF:tangential}.
\item The mean-field force $F^{(\alpha)}$ of~\eqref{eq:DF:mean_field_hilbert}
    becomes here $V^*\Lind(\rho)V$ in the core equation and
    $(I-VV^*)\Lind(\rho)V\sigma^{-1}$ in the frame equation.
\end{itemize}
The key difference is the additional constraint $\Tr(\dot\sigma) = 0$
(trace zero on the core), which has no analogue in the pure-state
MCTDH case.
\end{remark}

\subsection*{The closed kinematic system}

Combining~\eqref{eq:lind:core_explicit} and~\eqref{eq:lind:frame_eq},
the Dirac--Frenkel projected Lindblad dynamics in Stiefel-core
coordinates reads:
\begin{subequations}
\label{eq:lind:kinematic}
\begin{align}
    \dot\sigma(t)
    &= V(t)^*\, \Lind\bigl(V(t)\sigma(t)V(t)^*\bigr)\, V(t),
    \label{eq:lind:kinematic_core}\\[4pt]
    \dot{V}(t)
    &= \bigl(I - V(t)V(t)^*\bigr)\,
    \Lind\bigl(V(t)\sigma(t)V(t)^*\bigr)\,
    V(t)\, \sigma(t)^{-1},
    \label{eq:lind:kinematic_frame}
\end{align}
\end{subequations}
with initial conditions $V(0) = V_0 \in \Stiefel{r}{H}$ and
$\sigma(0) = V_0^* \rho_0 V_0 \in \Dpos{\CC^r}$.

\begin{proposition}[Well-posedness of the kinematic system]
\label{prop:lind:kinematic_wellposed}
Under Assumption~\ref{ass:lind:regular}, the system~\eqref{eq:lind:kinematic}
has a unique local solution $(V(t), \sigma(t))$ with
$V(t) \in \Stiefel{r}{H}$ and $\sigma(t) \in \Dpos{\CC^r}$ for
$t \in [0, T)$. The reconstructed density operator
$\rho(t) = V(t)\sigma(t)V(t)^*$ satisfies~\eqref{eq:lind:df}.
\end{proposition}

\begin{proof}
In the chart $\Phi_{S_0}$ centred at $\rho_0$, the system~\eqref{eq:lind:kinematic}
is a locally Lipschitz ODE on
$\cL(\CC^r, S_0^\perp) \times \cTsa{1}(\CC^r)$
(a Banach space, finite-dimensional in the core direction and infinite-
dimensional in the frame direction). Local existence and uniqueness
follow from Theorem~\ref{thm:app:picard}. Positivity and trace of $\sigma(t)$
are preserved as in Proposition~\ref{prop:lind:preservation}.
\end{proof}

\begin{remark}[Complexity and numerical implementation]
\label{rem:lind:complexity}
In a finite-dimensional truncation $H \cong \CC^d$ (dimension $d$,
rank $r \ll d$), the kinematic system~\eqref{eq:lind:kinematic}
requires storing only $V(t) \in \CC^{d \times r}$ and $\sigma(t) \in
\CC^{r \times r}$ at cost $O(dr + r^2)$, compared to $O(d^2)$ for the
full state $\rho(t)$. Each time step of~\eqref{eq:lind:kinematic}
costs $O(|J|dr^2 + dr^2 + r^3)$, where $|J|$ is the number of Lindblad
operators. The dominant cost is the evaluation of $\Lind(\rho)$
projected into the frame coordinates, which requires computing
$L_j V(t) \in \CC^{d \times r}$ for each $j$. This is a factor
$d/r$ reduction compared to the full simulation.
\end{remark}

%------------------------------------------------------------
\section{Historical notes}
\label{sec:lind:history}
%------------------------------------------------------------

\paragraph{The Lindblad equation.}
The GKSL master equation was derived independently by
Gorini, Kossakowski, and Sudarshan~(1976) and by Lindblad~(1976) as
the most general Markovian, completely positive, trace-preserving
quantum master equation. It is the quantum analogue of the Fokker--Planck
equation for classical Markov processes. The bounded-generator case
(Assumption~\ref{ass:lind:regular}) is completely understood;
the unbounded case remains an active research area.

\paragraph{Dirac--Frenkel for open systems.}
The application of the Dirac--Frenkel variational principle to the
Lindblad equation — projecting onto a finite-rank manifold — appears
in the numerical literature as the time-dependent variational principle
(TDVP) for density operators, used in the simulation of open quantum
spin chains via matrix product operators (MPO). The mathematical
foundations, including the well-posedness theorem
(Theorem~\ref{thm:lind:wellposed}) and the explicit bounded projector
(Proposition~\ref{prop:qs:projector}), are established in the companion
paper~\cite{AcevedoFalco2026}.

\paragraph{Low-rank methods for Lindblad equations.}
Numerical low-rank methods for the Lindblad equation were developed
by Le Bris, Rouchon, and Roussel~\cite{LeBrisRouchon2013,LeBrisRouchonRoussel2015}
in the Hilbert--Schmidt setting, and by Gravina and Savona~\cite{GravinaSavona2024}
using adaptive rank strategies. The CPTP-preserving approach via Kraus
parametrisations was developed by Appelö and Cheng~\cite{AppeloChengo2024}.
The present chapter provides a unified geometric framework for all
these methods.

\paragraph{Connection to tensor network methods.}
For multipartite open systems $H = \bigotimes_{j=1}^d H_j$, the density
operator $\rho(t)$ is a positive trace-class operator in
$\cTsa{1}(\bigotimes_j H_j)$. Approximating $\rho(t)$ by a matrix
product operator (MPO) of bond dimension $r$ is a special case of
the tree-based tensor format of Chapters~\ref{chap:topTBT}--\ref{chap:TB_manifold},
applied to the operator space $\cTsa{1}(\bigotimes_j H_j) \cong
\bigotimes_j \cL(H_j)$. The Dirac--Frenkel principle on this space
gives a rigorous foundation for TDVP-MPO algorithms for open systems.
This connection is explored in Open Problem~\ref{prob:lindblad_tensor}.
